\clearpage
\section{Anexos [Pendientes: 1. completar valoración y asignación de los registros de riesgos]}\label{sec:spmp-13}
Los anexos conservan registros complementarios que deben mantenerse durante el desarrollo, mientras que los procesos necesarios para comprenderlos permanecen descritos en el cuerpo del plan. En particular, el registro de riesgos y sus tratamientos se utilizan mediante la sección \ref{sec:spmp-12-2}. No se trasladan a este apartado la calendarización, los diagramas de procesos o la organización que la plantilla ubica en otras secciones~\cite{plantillaspmp}.
\subsection*{Registro inicial de riesgos [Pendientes: valorar probabilidad, impacto y prioridad]}\phantomsection\label{sec:riesgos-anexo}
Los identificadores siguientes no expresan un orden de importancia. La valoración y el orden por prioridad se completarán con la escala acordada, conservando fecha, responsable y evidencia de revisión. La tabla se actualizará cuando cambien las condiciones conocidas y no registrará como resuelto un riesgo únicamente porque se haya iniciado una gestión.
\begin{longtable}{|p{0.36\dimexpr\linewidth-34pt\relax}|p{0.39\dimexpr\linewidth-34pt\relax}|p{0.25\dimexpr\linewidth-34pt\relax}|}
\caption{Registro inicial de riesgos, pendiente de valoración. Fuente: brechas identificadas en el contexto del proyecto~\cite{datos2026,hpc2026,diseno2026}.}\\
\hline
\textbf{ID y riesgo} & \textbf{Efecto posible} & \textbf{Probabilidad, impacto y prioridad}\\\hline
\endfirsthead
\hline
\textbf{ID y riesgo} & \textbf{Efecto posible} & \textbf{Probabilidad, impacto y prioridad}\\\hline
\endhead
R01. Entradas insuficientes o no transferibles & Impiden configurar o justificar el caso territorial. & Por valorar.\\\hline
R02. Acceso o recursos HPC no confirmados & Retrasan ejecuciones o limitan generación y almacenamiento. & Por valorar.\\\hline
R03. Revisión WFDS o ejecución no reproducible & Impiden repetir casos y sustentar resultados. & Por valorar.\\\hline
R04. Agregación sin evidencia suficiente & Limita la interpretación de la aproximación y sus afirmaciones. & Por valorar.\\\hline
R05. Cobertura y costo de biblioteca inadecuados & Generan consultas sin representación suficiente o cómputo inviable. & Por valorar.\\\hline
R06. Datos meteorológicos no disponibles & Afectan las condiciones previstas del flujo de consulta. & Por valorar.\\\hline
R07. Alcance o aceptación sin acuerdo & Producen compromisos ambiguos y cambios tardíos. & Por valorar.\\\hline
R08. Interlocutores de evaluación no confirmados & Impiden realizar la evaluación prevista en sus condiciones. & Por valorar.\\\hline
R09. Indisponibilidad o pérdida de documentación & Afecta coordinación y recuperación de versiones. & Por valorar.\\\hline
R10. Costos y dedicación sin estimación suficiente & Afectan calendario, presupuesto y viabilidad. & Por valorar.\\\hline
\end{longtable}

\subsection*{Tratamientos iniciales [Pendientes: priorizar y asignar responsables, tiempo y recursos]}\phantomsection\label{sec:tratamientos-anexo}
Las acciones siguientes son propuestas vinculadas con los riesgos identificados y deberán priorizarse antes de comprometer recursos. Cada asignación nominal, fecha y esfuerzo se incorporará al registro del riesgo y a la actividad correspondiente del WBS, evitando que una lista general de respuestas sustituya un plan ejecutable.
\begin{longtable}{|p{0.15\dimexpr\linewidth-34pt\relax}|p{0.62\dimexpr\linewidth-34pt\relax}|p{0.23\dimexpr\linewidth-34pt\relax}|}
\caption{Prevención y tratamiento propuestos para ratificación. Fuente: elaboración propia con las dependencias del proyecto.}\\
\hline
\textbf{Riesgo} & \textbf{Prevención y tratamiento} & \textbf{Asignación}\\\hline
\endfirsthead
\hline
\textbf{Riesgo} & \textbf{Prevención y tratamiento} & \textbf{Asignación}\\\hline
\endhead
R01 y R03 & Revisar procedencia y entradas, conservar el caso y comprobar reproducción. Si falla, revisar configuración y limitar afirmaciones o alcance mediante decisión. & Responsable, esfuerzo y recursos por asignar.\\\hline
R02 y R05 & Confirmar condiciones y medir el piloto antes de comprometer generación. Si hay límites, reevaluar combinaciones y planificación con evidencia. & Coordinación con HPC y equipo por concretar.\\\hline
R04 & Diseñar casos de contraste y conservar discrepancias. Si la evidencia es insuficiente, explicitar límites y revisar el método antes de afirmar validez. & Revisión técnica y recursos por asignar.\\\hline
R06 & Definir y comprobar las condiciones de disponibilidad del servicio. Si no hay datos, aplicar la respuesta del producto que se acuerde en los requisitos. & Asignación dependiente de la elección de API.\\\hline
R07 y R08 & Conversar sobre aceptación, confirmar interlocutores y conservar acuerdos. Si cambia la disponibilidad, revisar las condiciones de evaluación y su impacto. & Coordinador e interlocutores por confirmar.\\\hline
R09 & Comprobar respaldo y recuperación y conservar versiones identificadas. Si ocurre pérdida, recuperar y contrastar consistencia de los artefactos. & Responsabilidad operativa por ratificar.\\\hline
R10 & Registrar esfuerzo, consultar tarifas y revisar estimaciones. Si aparecen desviaciones, ajustar recursos o compromisos mediante el proceso de cambios. & Coordinador y equipo por concretar.\\\hline
\end{longtable}
